%-------------------------------------------------------------------------------
%	SECTION TITLE
%-------------------------------------------------------------------------------
\cvsection{Özet}
\vspace{0mm}


%-------------------------------------------------------------------------------
%	CONTENT
%-------------------------------------------------------------------------------
\begin{cvparagraph}

%---------------------------------------------------------
Endüstriyel üretim sahasında dijital dönüşüm ve uygulamalı yapay zeka girişimlerine liderlik eden, 9 yılı aşkın deneyime sahip sonuç odaklı bir Dijitalleşme ve Yapay Zeka Ekip Lideri. Şu anda \textbf{3 üretim tesisinde} OT dijitalleşmesi ve yapay zeka stratejisini yürütüyor; 2 kişilik bir mühendislik ekibini yönetiyor, \textbf{10'dan fazla fonksiyon/departmanla} koordinasyon sağlıyor ve şirket çapındaki \textbf{AI Elçileri} programına liderlik ediyor. \textbf{Kestirimci bakım}, \textbf{bilgisayarlı görü tabanlı güvenlik sistemleri} (bu sistemlerin uygulandığı alanlarda \textbf{4 yıl kesintisiz sıfır iş kazası} sağlandı), \textbf{GenAI destekli otomasyon} ve \textbf{enerji \& sürdürülebilirlik analitiği} aracılığıyla ölçülebilir etki sağlıyor; saha verilerini sürekli olarak operasyonel verimliliğe, maliyet tasarrufuna ve daha güvenli çalışma koşullarına dönüştürüyor. \textbf{PLC}, \textbf{SCADA} ve endüstriyel protokoller (\textbf{OPC UA}, \textbf{MQTT}, \textbf{Profinet}) konusunda derin uygulamalı uzmanlığını, güçlü yazılım/bulut ve liderlik yetkinlikleriyle (\textbf{Node.js}, \textbf{React}, \textbf{Azure}/\textbf{AWS}, \textbf{Kubernetes}) dengeliyor. İnci Holding grup şirketleri çapında düzenlenen yarışmalarda kazanılan çoklu birincilik ve bir İnovasyon Özel Ödülü ile öne çıkıyor. Ayrıca uygulamalı makine öğrenmesi ve endüstriyel otomasyon alanında \textbf{4 farklı çalışması bilimsel yayın} olarak yayımlandı.
\end{cvparagraph}
